\documentclass[10pt]{article}

\usepackage{amsmath}
%
\usepackage{amssymb}
%
\usepackage{enumerate}
%
\usepackage[width=7in,height=10in,top=0.75in,centering]{geometry}
%
\usepackage[dvipsnames]{xcolor}
%
\usepackage{tikz}
%
\usepackage{multicol}
%
\usepackage{hyperref}
%
\usepackage{marginnote}
%
%%%%%%%%%%%%%%%% HEADER %%%%%%%%%%%%%%%%%%%%%
\usepackage{fancyhdr}
\pagestyle{fancy}
%%
%\lfoot{Left footer}
%\cfoot{Center of footer}
%\rfoot{Right of footer}
%%
\lhead{Survey of Calculus I}
\chead{Week 2 Workshop}
\rhead{\thepage}
%
%%%%%%%%%%%%%%%% HEADER %%%%%%%%%%%%%%%%%%%%%
%
%%%%%%%%%%%%%%%% PGFPLOTS %%%%%%%%%%%%%%%%%%%
\usepackage{pgfplots}
%
\pgfplotsset{every axis/.append style={
                    axis x line=middle,    % put the x axis in the middle
                    axis y line=middle,    % put the y axis in the middle
                    axis line style={<->,color=black}, % arrows on the axis
                    %xlabel={$x$},          % default put x on x-axis
                    ylabel={$y$},          % default put y on y-axis
            }}
%
 \pgfplotsset{compat=1.14}
%
 \usepgfplotslibrary{fillbetween}
%%%%%%%%%%%%%%%% PGFPLOTS %%%%%%%%%%%%%%%%%%%
%
\parindent0pt
\fboxsep=0.25cm
%
\newcommand{\Heading}[2]{\fbox{{\Large\textbf{#1}}}\hfill\fbox{\textbf{#2}}\par\vspace{0.1in}}
%
\newcommand{\Name}[1]{\textbf{Name}:\quad#1\hrulefill}
%
\newcommand{\Target}[1]{\textbf{Learning Target(s):}\quad\fbox{#1}\par\medskip}
%
\DeclareMathOperator{\arcsec}{arcsec}
\DeclareMathOperator{\arccot}{arccot}
\DeclareMathOperator{\arccsc}{arccsc}
%
\newcounter{example}[section]
\newcounter{exercise}[section]
%
\newcommand{\important}[2]{\begin{center}\fbox{\begin{minipage}{.97\textwidth}\noindent\textbf{#1}\par#2\end{minipage}}\end{center}}
%
\newcommand{\defn}[1]{\begin{center}\fbox{\begin{minipage}{.97\textwidth}\reversemarginpar\marginnote{\textbf{Definition}}#1\end{minipage}}\end{center}}
%
\newcommand{\example}[1]{\vspace{0.1in}{\textcolor{blue}{\textbf{Example \stepcounter{example}\theexample.}}}\hspace{0.1in}\reversemarginpar\marginnote{\fbox{\bf #1}}}
%
\newcommand{\exercise}[1]{\vspace{0.1in}{\textcolor{blue}{\textbf{Exercise \stepcounter{exercise}\Alph{exercise}.}}}\hspace{0.1in}\reversemarginpar\marginnote{\fbox{\bf #1}}}

%%%%%%%%%%%% BEGIN DOCUMENT %%%%%%%%%%%%%%%
%%%%%%%%%%%% BEGIN DOCUMENT %%%%%%%%%%%%%%%
%%%%%%%%%%%% BEGIN DOCUMENT %%%%%%%%%%%%%%%
%%%%%%%%%%%% BEGIN DOCUMENT %%%%%%%%%%%%%%%
%%%%%%%%%%%% BEGIN DOCUMENT %%%%%%%%%%%%%%%

\begin{document}

\thispagestyle{empty}

\Heading{MATH 1190 \& 1210 Survey of Calculus I}{\begin{minipage}{0.25\textwidth}\begin{center} Week 2 Workshop:\\ Simplifying Exponents \end{center}\end{minipage}}

\vspace{0.1in}

\Name{} \quad \Target{{\bf\color{RoyalBlue}L1}, {\bf\color{RoyalBlue}L2}}

\hrulefill\\

During week 3, we will compute limits involving expressions with exponents -- such as $x^{-1/2}$ and $e^{2x}$ -- and being comfortable manipulating these expressions will be helpful. So, let's get some practice rewriting expression with exponents!

\begin{enumerate}
\item[1.] Rewrite the following expressions with only positive exponents.\\
\end{enumerate}

\begin{align*}
x^{-3} &=\rule[-8pt]{1in}{0.4pt} & 2x^{-1} &=\rule[-8pt]{1in}{0.4pt} & \dfrac{1}{x^{-5}} &=\rule[-8pt]{1in}{0.4pt}\\[0.75in]
%
\dfrac{x^{-2}}{5} &=\rule[-8pt]{1in}{0.4pt} & \dfrac{2x^{-4}}{3} &=\rule[-8pt]{1in}{0.4pt} & \dfrac{4}{7x^{-6}} &=\rule[-8pt]{1in}{0.4pt}\\[0.75in]
%
6e^{-x} &=\rule[-8pt]{1in}{0.4pt} & \dfrac{1}{2e^{-x}} &=\rule[-8pt]{1in}{0.4pt} & \dfrac{e^{-x}}{x^2} &=\rule[-8pt]{1in}{0.4pt}\\[0.75in]
\end{align*}

\begin{enumerate}
\item[2.] Rewrite the following expressions using exponent notation. For example, $\dfrac{4}{\sqrt[3]{x^2}}$ would be rewritten as $4x^{-2/3}$.\\
\end{enumerate}

\begin{align*}
\sqrt[4]{x} &=\rule[-8pt]{1in}{0.4pt} & \sqrt{x^3} &=\rule[-8pt]{1in}{0.4pt} & \dfrac{1}{\sqrt[3]{x}} &=\rule[-8pt]{1in}{0.4pt}\\[0.75in]
%
\dfrac{\sqrt[3]{x+1}}{5} &=\rule[-8pt]{1in}{0.4pt} & \dfrac{5}{\sqrt[3]{x+1}} &=\rule[-8pt]{1in}{0.4pt} & \dfrac{2}{3\sqrt{x^2+1}} &=\rule[-8pt]{1in}{0.4pt}
\end{align*}

\newpage

For each group below separately, sketch a graph of a function satisfying all the given statements.\\

\begin{minipage}[t][2in]{0.45\textwidth}
Group: (A)
    \begin{itemize}
    \item $\displaystyle\lim_{x\to2}f(x) = 1$
    \item $f(2)=3$
    \item $\displaystyle\lim_{x\to4^+} f(x)=\infty$
    \item $\displaystyle\lim_{x\to4^-} f(x)=-\infty$
    \item $f(4)=1$\\
    \end{itemize}
\end{minipage}
%
\hfill
%
\begin{minipage}[t][2in]{0.45\textwidth}
Group: (a)
    \begin{itemize}
    \item $f$ has a removable discontinuity at $x=1$
    \item $f(1)=3$
    \item $f$ has an infinite discontinuity at $x=3$
    \item $f(3)=1$\\[0.2in]
    \end{itemize}
\end{minipage}
%
%
\vfill
%
%
\begin{minipage}[t][2in]{0.45\textwidth}
Group: (B)
    \begin{itemize}
    \item $\displaystyle\lim_{x\to2+}f(x)=1$
    \item $\displaystyle\lim_{x\to2^-}f(x)=3$
    \item $f(2)=2$
    \item $\displaystyle\lim_{x\to4}f(x)=2$
    \item $f(4)=4$\\
    \end{itemize}
\end{minipage}
%
\hfill
%
\begin{minipage}[t][2in]{0.45\textwidth}
Group: (b)
    \begin{itemize}
    \item $f$ has a jump discontinuity at $x=1$
    \item $f(1)=3$
    \item $f$ has a removable discontinuity at $x=3$
    \item $f(3)=4$\\[0.2in]
    \end{itemize}
\end{minipage}
%
%
\vfill
%
%
\begin{minipage}[t][2in]{0.45\textwidth}
Group: (C)
    \begin{itemize}
    \item $\displaystyle\lim_{x\to2^-}f(x)=2$
    \item $\displaystyle\lim_{x\to2^+}f(x)=3$
    \item $f(2)=2$
    \item $\displaystyle\lim_{x\to4}f(x)=\infty$
    \item $f(4)=3$\\
    \end{itemize}
\end{minipage}
%
\hfill
%
\begin{minipage}[t][2in]{0.45\textwidth}
Group: (c)
    \begin{itemize}
    \item $f(1)=2$
    \item $f$ has a jump discontinuity at $x=1$
    \item $f(3)=3$
    \item $f$ has an infinite discontinuity at $x=3$
    \end{itemize}
\end{minipage}
%
%
\vfill
%
%

%%%%%%%%%
\newpage%
%%%%%%%%%

What similarities do you notice between the graphs for group (A) and group (a), between group (B) and group (b), and between group (C) and group (c)?

\vfill

In terms of limits, what does it mean for $f$ to have a removable discontinuity at $x=a$?

\vfill

In terms of limits, what does it mean for $f$ to have a jump discontinuity at $x=a$?

\vfill

In terms of limits, what does it mean for $f$ to have an infinite discontinuity at $x=a$?

\vfill

\end{document}